% !TeX root = ../papers/nc-p3.tex
\section{階数二曲線を含む低次曲面と完全交叉の排除}
\label{sec:cofactor-low-surfaces}
引き続き被約四辺有限窓で一般可逆な全階数反対根を扱い、一般被約な余因子スキームの
正次元部分全体の被約台を $C$ とする。次数 $d$ と正規化上のPlücker次数 $\beta$ は前節と同じである。
\begin{theorem}\label{thm:cofactor-no-quadrics}
$C$ を含む二次曲面はない。$C$ が整ならその一般像二平面の族のPlücker次数は少なくとも三、
$h^1(\mathcal O_C(3))-h^1(\mathcal I_C(3))\ge2$ である。
\end{theorem}
\begin{proof}
任意の平面に含まれる $C$ の成分の総次数は四以下である。
総次数 $e\ge5$ なら、平面上のそれらの被約方程式 $f$ の平方が制限行列式を割る
（各成分の一般点で単元二行を消せばよい）。従って行列式は平面で恒等零であり、
階数二曲面の排除から一般階数三である。平面で各余因子は $f$ の倍数なので $e\le6$。
$f$ を除いた残りは次数一以下の階数一対称行列で、像は固定である。
実際多項式UFD上で階数一対称行列は $hvv^t$ ($v$ は原始列) であり、次数を比べればよい。
従って平面上で固定核列が消えて矛盾する。$C$ が二平面の和に載れば、両平面の被約曲線 $A,B$ に分け、
共通の直線成分があれば片方だけに含める。両方とも次数四以下で各三次切断は非特殊である。
両方が非空のとき、$A\cap B$ は交線上にあって長さ四以下
（交線を成分としない方と交線との交叉に含まれる）。
三次式はそこに全射するので貼合せ完全列より $H^1(\mathcal O_C(3))=0$。
一平面だけの場合も同じ消滅を得る。定理\ref{thm:cofactor-singular-special}に反する。

滑らかな二次曲面 $Q$ で $C$ の型を $(a,b)$ とする。曲面の完全列で
$h^1(\mathcal O_C(3))=h^2(\mathcal O_Q(3-a,3-b))>0$ だから $a,b\ge5$。
$(8,8)$ の行列式は $2C$ に沿って消えるので $Q$ 上で恒等零。
一般階数三だから各余因子から $C$ の方程式を除くと、非零の階数一対称行列で、
二次数は $(6-a,6-b)$、どちらも一以下である。二組の二項変数の多項式環で上と同じ
原始平方分解をすれば像は固定で、$Q$ 上で消える固定核列を得て矛盾する。

二次錐 $Q=\mathbb P(1,1,2)$ では、$C$ は重み付き多項式環で次数 $d$ の平方因子のない式 $f$ で切れる。
実際高さ一素因子を持つ斉次錐の和はこのUFDで主因子である。超平面は重み二であり、
一般超平面に制限すれば曲線の通常次数は $d$。
重み付き多項式環の中間コホモロジー消滅と標準加群のシフト $-4$ より
$h^1(\mathcal O_C(3))=h^0(\mathcal O_Q(d-10))$。ここでは $f$ 倍の次数別完全列を層化したので、
頂点で $C$ がCartierでない場合にも使える。従って $d\ge10$。
$2d>16$ より重み十六の行列式は錐上で消え、階数一の余因子を $f$ で除けば重み $12-d$ である。
$d\ge11$ ならその平方列は定数となり、やはり固定核列を生じる。
$d=10$ なら $f$ は五次超平面切断で、$C$ は二次と五次の完全交叉、$\chi(\mathcal O_C)=-15$。
命題\ref{prop:cofactor-singular-length}で $2m\ge60-30+\beta>30$。
しかし非零二次式上の印は高々十四（定理\ref{thm:reducible-octic-factors}の二次式の印数の証明は
階数二軌跡に条件を使わない）。これで全種類の二次曲面を尽くす。

整な場合、像の直線族は動く。その掃く曲面の次数と、一般像点を通る母線をパラメータ付きで
数えた数の積は $\beta$ である。実際正規化上の像二束の射影直線束で二つの一般出力超平面を引き戻すと、
両超平面の交線と母線が会う条件が一個のPlücker超平面を与える。
もし $\beta\le2$ なら全母線を含む非零の出力二次式 $q$ がある。$q$ は各像二空間上で零形式なので
$\operatorname{tr}(qN(x))$ は $C$ 上で零である。十成分の独立性からこれは非零の入力二次式で矛盾する。
最後の不等式は命題\ref{prop:cofactor-singular-length}による。
\end{proof}

\begin{proposition}\label{prop:cofactor-lci-kernel-extension}
$C$ が整な局所完全交叉なら一般核のGrassmann射は $C$ 全体には延長しない。
少なくとも一つの階数一特異点で $W$ は可逆でない。
\end{proposition}
\begin{proof}
延長があれば核二束 $U$ と $L=\det U^*$ を $C$ 上で取れる。六個の外積は $L$ を生成し、
$W\to L$ は全射かつ一般同型、捩れなしなので同型である。逆に $W$ が可逆なら
その生成切断による射は閉じたGrassmann多様体を経由する。
階数二の点では余核の二次外積は局所的に $\mathcal O_Z$ で、純商は $\mathcal O_C$ である。
滑らかな点では捩れなし階数一層は可逆である。

延長した $U$ の二切断を周囲に局所的に持ち上げて $N$ と両側から対にし、
$\mathcal I_C/\mathcal I_C^2$ に値を取れば、$NU=0$ により持ち上げの選択に依らない。
この形式の双対は $N_{C/\mathbb P^3}\to\operatorname{Sym}^2U^*(2)$ で、
単元二行の消去と一般被約性により一般階数二である。局所完全交叉の随伴と外積は
系\ref{cor:all-smooth-rank-two-excess-exclusion}と同じく $\omega_C\to (L^3)^{\oplus10}$ を与え、
従って $2p_a-2\le3\beta$。二次外積行列のスカラー部分は $\mathcal O_C(4)\otimes L^{-2}$
の非零切断で、$m\le4d-2\beta$。各相異なる印はその零スキームに少なくとも長さ一で寄与する。
前定理による特殊性から $2p_a-2\ge3d$、よって $\beta\ge d$。
命題\ref{prop:cofactor-singular-length}で
$0\le2m-6d+2p_a-2-\beta\le2(d-\beta)\le0$。
等号なら $2p_a-2=3d$、次数零の線束 $\omega_C(-3)$ は非零切断を持つので自明となり
$h^1(\mathcal O_C(3))=1$。前定理に反する。
\end{proof}

\begin{theorem}\label{thm:cofactor-no-integral-ci}
$C$ は整な大域完全交叉ではあり得ない。
\end{theorem}
\begin{proof}
生成次数を $a\le b$ とし、方程式を $f,g$ とする。前定理より $a\ge3$。
余因子の共通零に曲面はないので $b\le6$。$\det N\in(f,g)^2$ である。
実際 $C$ の一般点で二行消去により二乗に入り、完全交叉イデアルの二乗には埋込素点がない
（$(f,g)/(f,g)^2$ は $C$ 上自由）。特に $a\le4$。
$b>a$ なら $f$ は既約であり、$a=b$ なら生成鉛筆のどの元も既約である。
可約ならその一因子が整な $C$ を含み、最小次数に反するからである。
特殊性と随伴より $a+b\ge7$。

$a=3,b=4$ なら $\omega_C=\mathcal O_C(3)$ が上の下界 $h^1\ge2$ に反する。
$b\ge5$ なら $8<2b$ なので $f\mid\det N$。$F=V(f)$ の一般階数は三であり、
その上で余因子行列は $gB$ で $B$ の成分の次数は $6-b\le1$。
$B$ の二行小行列式の次数は二以下で $F$ 上零なので恒等零、原始平方分解から像が固定となる。
すると $F$ で固定核列が消え、四つの独立な二次式が同時に $f$ で割れることになり矛盾する。

残る $(4,4)$ では $\det N$ は $f,g$ の定数係数二次形式なので四次の鉛筆の二元の積である。
一方を新たな $f$ に取り、別の独立な生成元を $g$ とする。$F$ 上で
$\operatorname{adj}N=gB$、$B$ は周囲の対称二次式行列であり、$\bigwedge^2 B=fD$、$D$ は定数行列である。
$D=0$ なら $B=hvv^t$、$\deg v\le1$、ゆえに $Nv$ の全成分が $F$ 上零で次数三以下、
恒等零となり一般可逆性に反する。
$D\ne0$ とし、一般階数を $r\ge2$ とする。$B$ の像二次外積の像は固定なので
$B$ の像 $V$ 自体が固定される（$\bigwedge^2 V$ の全収縮像が $V$ を張る）。
$r\ge3$ なら $V$ 上の可逆な対称 $r$ 行部分の二次外積は射影的に一定であり、
元の可逆写像も射影的に一定である。実際二次外積が射影的に恒等な変換は
各二空間を保ち、その交叉で各一点を保つ。従って $B$ は定数行列の二次式倍となり、
二次式の平方が既約 $f$ の定数倍とはなれない。
$r=2$ なら $F$ 上の $N$ の核点は固定直線 $\mathbb P(V)$ に入る。
この二空間の基底 $v,w$ に対する $Nv\wedge Nw$ は四次で $F$ 上零なので $f$ と定数外積の積。
一般可逆性から非零で、その一般像二平面は固定である。
特に $Nv$ の四つの二次式が独立であることに反する。これで全生成次数を除いた。
\end{proof}
